\begin{table}[ht]
\centering
\small
\caption{Improvement factor over the brute-force baseline as the precision requirement tightens, at $\phi\sim\mathcal{U}(0.01,0.1)$ and $90\%$ convergence. The advantage grows with precision and then saturates; the ceiling column bounds how much of it is attainable at all.}
\label{tab:scaling-with-prec}
\begin{tabular}{lcccc}
\toprule
Precision $\epsilon$
& LS
& BS
& Rev. Eng.
& \textit{Ceiling} \\
\midrule
$10^{-3}$
& 1.08\texttimes
& 1.42\texttimes
& 1.54\texttimes
& 1.86\texttimes \\
$10^{-4}$
& 1.56\texttimes
& 1.60\texttimes
& 1.74\texttimes
& 1.83\texttimes \\
$10^{-5}$
& 1.57\texttimes
& 1.75\texttimes
& 1.80\texttimes
& 1.82\texttimes \\
$10^{-6}$
& 1.63\texttimes
& 1.78\texttimes
& 1.80\texttimes
& 1.83\texttimes \\
$10^{-7}$
& 1.65\texttimes
& 1.79\texttimes
& 1.81\texttimes
& 1.83\texttimes \\
$10^{-8}$
& 1.62\texttimes
& 1.77\texttimes
& 1.82\texttimes
& 1.83\texttimes \\
\bottomrule
\end{tabular}
\\[2pt]
\footnotesize Adaptive entries are de-biased: parameters are grid-tuned on seeds 42 and 43 and the frozen winner is re-validated on the independent seed 2024 at $R=50{,}000$ trials. The ceiling is exact -- its convergence probability is known in closed form for each $\phi$ and averaged over the prior -- so it carries no Monte-Carlo error and no interval.
\end{table}
